% Part: first-order-logic
% Chapter: completeness
% Section: completeness-theorem

\documentclass[../../../include/open-logic-section]{subfiles}

\begin{document}

\iftag{FOL}
      {\olfileid{fol}{com}{cth}}
      {\olfileid{pl}{com}{cth}}

\olsection{পূর্ণতা উপপাদ্য}

\begin{explain}
আমাদের ফলগুলি একত্র করি: তাতে পূর্ণতা উপপাদ্যে পৌঁছাই।
\end{explain}

\begin{thm}[পূর্ণতা উপপাদ্য]
\ollabel{thm:completeness}
$\Gamma$-কে !!{sentence}s-এর একটি সেট ধরি। $\Gamma$ সঙ্গত হলে সেটি
পরিতৃপ্তিযোগ্য।
\end{thm}

\begin{proof}
ধরা যাক $\Gamma$ সঙ্গত।\iftag{FOL}{ \olref[hen]{lem:henkin} অনুসারে
  একটি সম্পৃক্ত সঙ্গত সেট $\Gamma' \supseteq \Gamma$ আছে।}{}
\olref[lin]{lem:lindenbaum} অনুসারে এমন একটি $\Gamma^* \supseteq
\iftag{FOL}{\Gamma'}{\Gamma}$ আছে, যা সঙ্গত ও !!{complete}।
\iftag{FOL}{
  $\Gamma' \subseteq \Gamma^*$ বলে প্রতিটি !!{formula}~$!A(x)$-এর জন্য
  $\Gamma^*$-তে নিচের আকারের একটি !!a{sentence} আছে:
  \iftag{prvEx}
        {$\lexists[x][!A(x)] \lif !A(c)$}
        {$\lnot\lforall[x][!A(x)] \lif \lnot !A(c)$}।
  তাই $\Gamma^*$ সম্পৃক্ত। $\Gamma$-তে~$\eq$ না থাকলে}{তখন}
\olref[mod]{lem:truth} অনুসারে
$\iftag{FOL}{\Sat{M(\Gamma^*)}}{\pSat{v(\Gamma^*)}}{!A}$ যদি এবং কেবল
যদি $!A \in \Gamma^*$। বিশেষত, এখান থেকে প্রতিটি $!A \in \Gamma$-র
জন্য $\iftag{FOL}{\Sat{M(\Gamma^*)}}{\pSat{v(\Gamma^*)}}{!A}$ পাওয়া
যায়; তাই $\Gamma$ পরিতৃপ্তিযোগ্য।\iftag{FOL}{
  $\Gamma$-তে~$\eq$ থাকলে \olref[ide]{lem:truth} অনুসারে সব
  !!{sentence}s~$!A$-এর জন্য
  $\Sat{\equivclass{M}{\approx}}{!A}$ যদি এবং কেবল যদি $!A \in
  \Gamma^*$। বিশেষত, প্রতিটি $!A \in \Gamma$-র জন্য
  $\Sat{\equivclass{M}{\approx}}{!A}$; তাই $\Gamma$ পরিতৃপ্তিযোগ্য।}{}
\end{proof}

\begin{cor}[পূর্ণতা উপপাদ্য, দ্বিতীয় রূপ]
\ollabel{cor:completeness}
প্রতিটি $\Gamma$ এবং !!{sentence}s~$!A$-এর জন্য: $\Gamma \Entails !A$
হলে $\Gamma \Proves !A$।
\end{cor}

\begin{proof}
লক্ষ করো, \olref{cor:completeness} ও \olref{thm:completeness}-এর
$\Gamma$-গুলি সার্বিকভাবে পরিমাণিত। যেন বিভ্রান্ত না হই, তাই অন্য একটি
চল ব্যবহার করে \olref{thm:completeness}-কে আবার বলি: !!{sentence}s-এর
যেকোনো সেট~$\Delta$-র জন্য, $\Delta$ সঙ্গত হলে সেটি পরিতৃপ্তিযোগ্য।
প্রতিবিপরীতকরণ অনুসারে, $\Delta$ পরিতৃপ্তিযোগ্য না হলে $\Delta$ অসঙ্গত।
উপসিদ্ধান্তটি প্রমাণ করতে এই কথাটি ব্যবহার করব।

ধরা যাক $\Gamma \Entails !A$। তাহলে \olref[syn][sem]{prop:entails-unsat}
অনুসারে $\Gamma \cup \{\lnot !A\}$ অপরিতৃপ্তিযোগ্য। আমাদের~$\Delta$
হিসেবে $\Gamma \cup \{\lnot !A\}$ নিলে, \olref{thm:completeness}-এর
আগের রূপ থেকে পাই যে $\Gamma \cup \{\lnot !A\}$ অসঙ্গত। নিচের ফলগুলি
অনুসারে $\Gamma \Proves !A$:
\iftag{FOL}{%
  \tagrefs{prfAX/{fol:axd:prv:prop:prov-incons},
    prfSC/{fol:seq:prv:prop:prov-incons},
    prfND/{fol:ntd:prv:prop:prov-incons},
    prfTab/{fol:tab:prv:prop:prov-incons}}}{%
  \tagrefs{prfAX/{pl:axd:prv:prop:prov-incons},
    prfSC/{pl:seq:prv:prop:prov-incons},
    prfND/{pl:ntd:prv:prop:prov-incons},
    prfTab/{pl:tab:prv:prop:prov-incons}}}।
\end{proof}

\tagprob{FOL}
\begin{prob}
\olref[fol][com][cth]{cor:completeness} ব্যবহার করে
\olref[fol][com][cth]{thm:completeness} প্রমাণ করো; এভাবে দেখাও যে
পূর্ণতা উপপাদ্যের দুটি রূপ সমতুল্য।
\end{prob}
\tagendprob

\tagprob{notFOL}
\begin{prob}
\olref[pl][com][cth]{cor:completeness} ব্যবহার করে
\olref[pl][com][cth]{thm:completeness} প্রমাণ করো; এভাবে দেখাও যে
পূর্ণতা উপপাদ্যের দুটি রূপ সমতুল্য।
\end{prob}
\tagendprob

\tagprob{FOL}
\begin{prob}
কোনো !!a{derivation} পদ্ধতি পূর্ণ হতে হলে তার বিধিগুলিকে এমন শক্তিশালী
হতে হবে, যাতে প্রতিটি অপরিতৃপ্তিযোগ্য সেটকে অসঙ্গত প্রমাণ করা যায়। পূর্ণতা
প্রমাণ করতে !!{derivation}-এর কোন বিধিগুলি প্রয়োজন ছিল? এই বিধিগুলির
কোনোটি কি প্রমাণের কোথাও ব্যবহৃত হয়নি? প্রশ্নগুলির উত্তর দিতে এমন একটি
তালিকা বা চিত্র তৈরি করো, যেখানে
\olref[fol][com][cth]{thm:completeness}-এর প্রমাণের আগের কোন ফলে
!!{derivation}-এর কোন বিধি ব্যবহৃত হয়েছে তা দেখানো থাকে। এই প্রমাণগুলিতে
বিধির যেসব অনুক্ত ব্যবহার আছে, সেগুলিও অবশ্যই উল্লেখ করো।
\end{prob}
\tagendprob

\tagprob{notFOL}
\begin{prob}
কোনো !!a{derivation} পদ্ধতি পূর্ণ হতে হলে তার বিধিগুলিকে এমন শক্তিশালী
হতে হবে, যাতে প্রতিটি অপরিতৃপ্তিযোগ্য সেটকে অসঙ্গত প্রমাণ করা যায়। পূর্ণতা
প্রমাণ করতে !!{derivation}-এর কোন বিধিগুলি প্রয়োজন ছিল? এই বিধিগুলির
কোনোটি কি প্রমাণের কোথাও ব্যবহৃত হয়নি? প্রশ্নগুলির উত্তর দিতে এমন একটি
তালিকা বা চিত্র তৈরি করো, যেখানে
\olref[pl][com][cth]{thm:completeness}-এর প্রমাণের আগের কোন ফলে
!!{derivation}-এর কোন বিধি ব্যবহৃত হয়েছে তা দেখানো থাকে। এই প্রমাণগুলিতে
বিধির যেসব অনুক্ত ব্যবহার আছে, সেগুলিও অবশ্যই উল্লেখ করো।
\end{prob}
\tagendprob

\end{document}
